% chapters/ch01_embedded_c.tex — Chapter 1: Embedded C Mastery
% Part 1: pointers -> arrays -> strings -> structs/unions/bitfields -> padding.
% Part 2 (volatile/const/static/storage classes/typecasting/preprocessor/
% endianness) will be APPENDED to this file in a later session.
% All compilable snippets verified with: gcc -Wall -Wextra -std=c11 (zero warnings).

\chapter{Embedded C Mastery}
\label{ch:embedded-c}

\begin{partnote}
\textbf{Part 1 of 2} --- Pointers $\rightarrow$ Arrays $\rightarrow$ Strings
$\rightarrow$ Structs / Unions / Bitfields $\rightarrow$ Padding \& Alignment.
Part 2 (appended to this chapter later) covers \code{volatile}, \code{const},
\code{static}, storage classes, typecasting, the preprocessor, and endianness.
\end{partnote}

This is the make-or-break chapter. In an embedded or avionics interview, every
panel eventually returns to C fundamentals, because the language \emph{is} the
hardware interface. The questions below are the ones actually asked at
screening, panel, and principal-engineer level --- answer them out loud until
they are reflexes.

%=====================================================================
\section{Primer}
%=====================================================================

\subsection{What a pointer actually is}

A pointer is a variable whose value is a memory address. On a 32-bit target
(e.g.\ the LPC3250/ARM9 in the BEL work, or a Cortex-M autopilot board) a
pointer is 4 bytes regardless of what it points to; on a 64-bit host it is
8 bytes. The \emph{type} of a pointer does not change its size --- it tells the
compiler two things: how many bytes to read or write when the pointer is
dereferenced, and how far to move when the pointer is incremented.

\begin{cgym}
int x = 42;
int *p = &x;   /* p holds the address of x    */
*p = 7;        /* dereference: write 7 into x */
\end{cgym}

\code{\&} takes an address; \code{*} (in an expression) dereferences. In a
\emph{declaration}, \code{*} is part of the type: \code{int *p} means
``\code{p} is a pointer to int.''

\subsection{Pointer arithmetic is typed}

\code{p + 1} does \textbf{not} add 1 byte. It adds
\code{1 * sizeof(*p)} bytes. This is the single most common source of
confusion and the root of many off-by-$N$ memory bugs.

\begin{cgym}
int arr[4] = {10, 20, 30, 40};
int *p = arr;     /* points at arr[0]                       */
p = p + 2;        /* advances 2*4 = 8 bytes on a 32-bit int */
/* *p is now 30 */
\end{cgym}

The difference of two pointers into the same array yields the number of
\emph{elements} between them (type \code{ptrdiff\_t}), not bytes.

\subsection{The hardware reason this matters}

Embedded code constantly casts an integer address to a pointer to talk to a
peripheral register:

\begin{cgym}
/* target-only: memory-mapped register access */
#define UART0_DR (*(volatile unsigned int *)0x40011004u)
UART0_DR = 'A';   /* write a byte to the UART data register */
\end{cgym}

If the pointer type is wrong, the access width is wrong, and a 32-bit
peripheral register read as 8 bits will silently corrupt your driver. This is
exactly the class of low-level peripheral-driver verification done on the BEL
UART/GPIO drivers.

\subsection{Arrays are not pointers (but they decay)}

An array name is \textbf{not} a pointer. It is a block of contiguous storage.
In \emph{most} expressions the array name ``decays'' to a pointer to its first
element --- but not under \code{sizeof}, not under \code{\&}, and not as a
string-literal initialiser. Confusing the two is a Tier-1 elimination
question.

\subsection{Structs, unions, bitfields, padding}

\begin{itemize}
\item \textbf{struct}: members laid out in order, each at its own offset;
  total size is often larger than the sum of members because the compiler
  inserts \emph{padding} to satisfy alignment.
\item \textbf{union}: all members share the same storage; size is that of the
  largest member. Used for type-punning and for memory-tight
  register/overlay tricks.
\item \textbf{bitfield}: members specified in bits, packed into words --- used
  for register maps and protocol headers, but with implementation-defined
  ordering you must know before trusting one on the wire.
\end{itemize}

These primer pages are the minimum. The depth lives in the Q\&A and the
Coding Gym below.

%=====================================================================
\section{Q\&A Bank}
%=====================================================================

\tierbanner{Tier 1 --- Screening (phone / HR-technical)}%
{Answer in 2--4 sentences. These are elimination questions: a hesitation here
ends the call.}

\question{Q1.1}{How big is a pointer?}
It is the size of an address on the target, not the size of what it points
to. On a 32-bit MCU every pointer is 4 bytes; on a 64-bit host, 8 bytes. So
\code{sizeof(char*) == sizeof(double*)} on any given platform.

\question{Q1.2}{What is the difference between \code{*p} in a declaration and
\code{*p} in an expression?}
In a declaration, \code{int *p} declares \code{p} as a pointer to int --- the
\code{*} binds to the type. In an expression, \code{*p} dereferences \code{p}
to read or write the pointed-to object. Same symbol, two roles.

\question{Q1.3}{What does \code{p++} do when \code{p} is an \code{int*}?}
It advances \code{p} by \code{sizeof(int)} bytes (4 on most targets), so it
now points at the next \code{int}, not the next byte. Pointer arithmetic is
always scaled by the pointed-to type.

\question{Q1.4}{What is a NULL pointer and why use it?}
\code{NULL} is a pointer value guaranteed not to point at any object. It is
used as a sentinel --- ``no target yet'' or ``end of list'' --- and
dereferencing it is undefined behaviour, which on an MCU usually means a fault
or a bogus read of the vector table at address 0.

\question{Q1.5}{Difference between an array and a pointer in one line?}
An array is a fixed block of storage; a pointer is a variable holding an
address. The array name decays to a pointer in most expressions, but
\code{sizeof} of an array gives the whole block while \code{sizeof} of a
pointer gives the address size.

\question{Q1.6}{What terminates a C string?}
A single NUL byte, \code{'\textbackslash{}0'} (value 0). Every standard string
function relies on it; a buffer without it is not a string and \code{strlen}
will run off the end.

\question{Q1.7}{Is \code{sizeof} a function or an operator?}
An operator, evaluated at compile time (for non-VLA types). It yields
\code{size\_t}. Because it is compile-time,
\code{sizeof(arr)/sizeof(arr[0])} gives the element count of a real array ---
but only where the array has not decayed to a pointer.

\question{Q1.8}{Why is a \code{struct} sometimes bigger than the sum of its
members?}
Because the compiler inserts padding bytes so each member sits at a properly
aligned address, and pads the tail so arrays of the struct stay aligned.
Reordering members largest-to-smallest usually shrinks the struct.

\question{Q1.9}{What is a dangling pointer?}
A pointer that still holds the address of memory that is no longer valid ---
a freed heap block, or a local variable whose function has returned.
Dereferencing it is undefined behaviour; the classic embedded case is
returning the address of a stack buffer from a function.

\question{Q1.10}{What is a \code{void} pointer?}
A typeless address used for generic interfaces --- \code{memcpy},
\code{qsort}, callback context arguments. You cannot dereference it or do
arithmetic on it in standard C; you must first cast it to a concrete pointer
type that says how many bytes an access means.

\tierbanner{Tier 2 --- Core technical (panel round)}%
{Answer in a short paragraph. The panel is checking that you have used these,
not just read about them.}

\question{Q2.1}{Explain pointer-to-pointer and give a real use.}
A pointer-to-pointer (\code{int **pp}) holds the address of a pointer. The
classic use is a function that must modify the caller's pointer --- for
example an allocator or a linked-list \code{push} that needs to update the
caller's \code{head}. C passes everything by value, so to change a
\code{Node*} in the caller you pass its address as \code{Node**}. Another
everyday use is \code{argv} (\code{char **argv}) and ragged 2D arrays where
each row is a separately allocated buffer.

\question{Q2.2}{Walk through \code{const} placement: \code{const char *p} vs
\code{char *const p} vs \code{const char *const p}.}
Read the declaration right-to-left from the variable name.
\code{const char *p} --- ``\code{p} is a pointer to char that is const'': you
can repoint \code{p}, but you cannot write through it (\code{*p = 'x'} is
illegal). \code{char *const p} --- ``\code{p} is a const pointer to char'':
you can write \code{*p}, but you cannot repoint \code{p}.
\code{const char *const p} --- both fixed. In drivers, a read-only register
view is often \code{volatile const uint32\_t *}, and string-literal
parameters should be \code{const char *} so the compiler stops you writing
into read-only memory.

\question{Q2.3}{What is array-to-pointer decay, and where does it NOT happen?}
In nearly every expression an array name converts to a pointer to its first
element --- passing to a function, arithmetic, comparison. Decay does
\textbf{not} happen under \code{sizeof} (you get the full array size), under
unary \code{\&} (\code{\&arr} has type ``pointer to array'', a different
type), and for a string literal used to initialise a \code{char[]}. The
practical trap: pass an array to a function and \code{sizeof} inside the
function gives the pointer size, so you must pass the length separately.

\question{Q2.4}{How do you declare and use a function pointer? Why does
embedded code use them?}
\code{int (*fp)(int, int);} declares \code{fp} as a pointer to a function
taking two ints and returning int; assign with \code{fp = \&add;} (or just
\code{add}) and call with \code{fp(2, 3)}. Embedded firmware uses them for
jump/dispatch tables (command handlers, protocol-message parsers), for
callback registration (ISR-to-application hooks), and for the interrupt
vector table itself, which is an array of function pointers. They replace
long \code{switch} ladders with O(1) dispatch and make state machines
table-driven.

\question{Q2.5}{Explain a \code{union} and one legitimate embedded use.}
A union overlays all its members in the same storage, so its size equals the
largest member and writing one member overwrites the others. Legitimate uses:
(1) type-punning to inspect the raw bytes of a value --- e.g.\ checking
endianness or serialising a float into a telemetry packet; (2) saving RAM
when several mutually-exclusive data shapes are never live at once;
(3) register overlays where the same word is sometimes read as a whole and
sometimes as named bitfields. The discipline is that you must track which
member is currently valid yourself --- the union does not.

\question{Q2.6}{What is structure padding and how do you control it?}
The compiler aligns each member to its natural boundary (a \code{uint32\_t}
to a 4-byte address, etc.)\ by inserting padding, and pads the struct tail to
the largest member's alignment. You control it by ordering members
largest-to-smallest to minimise gaps, or --- when a struct must match a
hardware register layout or a wire format exactly --- by packing it
(\code{\#pragma pack} or \code{\_\_attribute\_\_((packed))}). Packing removes
padding but can create unaligned members, which on some cores (older ARM)
faults or is slow, so it is used deliberately, not by default.

\question{Q2.7}{Why must protocol/packet structs be handled carefully across
machines?}
Two reasons collide: padding and endianness. A struct mapped onto a received
byte buffer may have different padding than the sender assumed, and
multi-byte fields may be byte-swapped if the two ends differ in endianness.
The robust approach for something like a MIL-STD-1553B or RS422 telemetry
frame is to parse byte-by-byte (or use a \code{packed} struct \emph{and}
explicit byte-order conversion), never to blindly cast the raw buffer to a
normal struct pointer. This is exactly why the BEL packet-parsing automation
worked at the byte level.

\question{Q2.8}{Why is dynamic memory (\code{malloc}) avoided in flight
software?}
Three reasons. \emph{Determinism}: \code{malloc}'s execution time varies with
heap state, which breaks worst-case timing analysis in a hard real-time loop.
\emph{Fragmentation}: over a long mission the heap fragments and an
allocation can fail even with enough total free memory --- and there is
rarely a sane way to ``handle'' that mid-flight. \emph{Certification}:
DO-178B-style verification wants memory usage known and bounded at build
time, so flight code allocates everything statically or from fixed-size
pools created once at initialisation (MISRA C likewise bans dynamic
allocation after init). The interview-ready summary: ``allocate at init,
never in the control loop.''

\question{Q2.9}{How does \code{qsort} use \code{void*} and function pointers
together?}
\code{qsort(base, n, size, cmp)} sorts an array of \emph{anything}: it takes
the array as \code{void*}, the element size, and a comparator function
pointer \code{int (*cmp)(const void*, const void*)}. Inside the comparator
you cast both \code{const void*} arguments back to the real element type and
return negative/zero/positive. This is C's idiom for generic programming ---
the same pattern appears all over firmware as driver callbacks that receive a
\code{void *ctx} context pointer, letting one driver serve many clients
without knowing their types.

\tierbanner{Tier 3 --- Trap \& depth (principal-engineer level)}%
{A paragraph, plus the trap the question defends against. These decide
senior-level offers.}

\question{Q3.1}{Is \code{arr[i]} the same as \code{*(arr + i)}? Is
\code{i[arr]} legal?}
Yes --- \code{arr[i]} is \emph{defined} as \code{*(arr + i)}. Because
addition is commutative, \code{*(arr + i) == *(i + arr)}, so \code{i[arr]} is
legal C and compiles to the same access. \trap interviewers use
\code{i[arr]} to test whether you understand that subscripting is pure
pointer arithmetic, not a special array operation. The follow-up trap is that
this equivalence is why there is no bounds checking in C ---
\code{arr[i]} will happily compute an out-of-range address, which is the
origin of buffer-overrun defects.

\question{Q3.2}{What is the difference between \code{char *s = "hi";} and
\code{char s[] = "hi";}?}
\code{char s[] = "hi"} copies the 3 bytes (\code{h}, \code{i},
\code{\textbackslash{}0}) into a writable local array you own --- you may
modify \code{s[0]}. \code{char *s = "hi"} makes \code{s} point at a string
\emph{literal}, which lives in read-only storage; writing \code{s[0] = 'H'}
is undefined behaviour and on an MCU typically faults or is silently dropped
if the literal is in flash. \trap the second form should really be
\code{const char *s} and modern compilers warn; candidates who say ``they're
the same'' fail. Also \code{sizeof(s)} is 3 for the array but the pointer
size for the pointer.

\question{Q3.3}{Given \code{int a[3][4];}, what are the types of \code{a},
\code{a[0]}, \code{*a}, \code{\&a}, and how does \code{a+1} differ from
\code{a[0]+1}?}
\code{a} has type ``array[3] of array[4] of int''. In an expression it decays
to \code{int (*)[4]} --- pointer to a row of 4 ints. \code{a[0]} is
``array[4] of int'', decaying to \code{int*}. \code{*a} is the same as
\code{a[0]}. \code{\&a} is \code{int (*)[3][4]} --- pointer to the whole 2D
array. The killer: \code{a + 1} advances by one \emph{row} =
\code{4*sizeof(int)} = 16 bytes, while \code{a[0] + 1} advances by one
\emph{int} = 4 bytes. \trap this is how interviewers check you understand
multidimensional arrays are not arrays of pointers --- they are one
contiguous block, and the row-pointer type carries the inner dimension. Pass
such an array as \code{void f(int (*a)[4], int rows)}, never as
\code{int **}.

\question{Q3.4}{With \code{union \{ float f; uint32\_t u; \} x;}, you set
\code{x.f} and read \code{x.u} --- legal? What about
\code{*(uint32\_t*)\&f}?}
Reading a different union member than the one last written is
\textbf{type-punning}. In C (since C99, with a footnote) reading through a
union member is permitted and is the portable, defined way to reinterpret
bytes --- the right tool for inspecting a float's IEEE-754 bits or
serialising it. The cast form \code{*(uint32\_t*)\&f} instead violates the
\textbf{strict-aliasing} rule: the compiler may assume a \code{float*} and a
\code{uint32\_t*} never refer to the same object, so with optimisation it can
reorder or elide the access and produce wrong code. \trap many engineers
reach for the pointer cast; the correct, optimiser-safe answers are the union
or \code{memcpy}. GCC's strict aliasing (on at \code{-O2}) will bite the cast
version.

\question{Q3.5}{A struct has a \code{uint8\_t}, then a \code{uint32\_t}, then
a \code{uint16\_t}. Sketch the layout, give its size, and show how to halve
the padding.}
Natural alignment on a 32-bit target: \code{uint8\_t} at offset 0, then
3 padding bytes, \code{uint32\_t} at offset 4, \code{uint16\_t} at offset 8,
then 2 tail-padding bytes so the struct size (12) is a multiple of its
largest alignment (4). Size = 12, of which 5 bytes are padding. Reorder to
\code{uint32\_t, uint16\_t, uint8\_t}: offsets 0, 4, 6, then 1 tail byte
$\rightarrow$ size 8, only 1 padding byte. \trap the interviewer wants the
\emph{tail} padding explained (so arrays of the struct stay aligned) and
wants you to know reordering --- not packing --- is the first, free fix;
packing is reserved for wire/register layouts because it can introduce
unaligned accesses.

\question{Q3.6}{Bitfields: name three things that are implementation-defined,
and why that disqualifies them from defining a wire protocol.}
(1) The \emph{allocation order} of bits within a unit (MSB-first vs
LSB-first) is implementation-defined; (2) whether a bitfield may straddle a
storage-unit boundary; (3) the signedness of a plain \code{int} bitfield
(\code{int x:1} may hold values 0 and $-1$). There is also padding between
bitfields and alignment of the underlying unit. \trap because two compilers
(or the same compiler on two architectures) can lay the same bitfield struct
out differently, a bitfield struct cast onto a received CCDL/1553B/RS422
frame may parse correctly on one build and garble on another. For
on-the-wire or cross-channel data you extract fields with explicit shifts and
masks; bitfields are fine only for local, single-compiler register
convenience.

\question{Q3.7}{Explain \code{restrict} and one place it earns its keep in a
driver.}
\code{restrict} on a pointer parameter is a promise to the compiler that, for
the lifetime of that pointer, the object it points to is accessed \emph{only}
through that pointer (no aliasing). It lets the optimiser keep values in
registers instead of reloading after every store. A \code{memcpy}-style
routine declared
\code{void copy(uint8\_t *restrict dst, const uint8\_t *restrict src, size\_t n)}
can be vectorised/unrolled because the compiler knows \code{dst} and
\code{src} don't overlap. \trap if you pass overlapping buffers to a
\code{restrict} function the behaviour is undefined --- which is precisely
why \code{memcpy} is \code{restrict} and \code{memmove} is not. Misusing
\code{restrict} to silence a warning is a real-world bug source.

\question{Q3.8}{What is a flexible array member, and when would you use one
for telemetry?}
C99 lets the \emph{last} member of a struct be an incomplete array:
\code{struct msg \{ uint16\_t id; uint16\_t len; uint8\_t payload[]; \};}.
\code{sizeof(struct msg)} excludes the payload; you allocate
\code{sizeof(struct msg) + len} bytes and the payload sits contiguously after
the header --- a natural fit for variable-length telemetry or CCDL frames
where a fixed header announces the payload size. \trap three of them: a
struct with a flexible array member cannot be an array element or another
struct's middle member; \code{sizeof} silently excludes the payload (a
classic buffer-size bug); and the pre-C99 \code{payload[0]} or
\code{payload[1]} ``struct hack'' is not portable --- say ``flexible array
member, C99'' and you signal real standard knowledge.

%=====================================================================
\section{Coding Gym}
%=====================================================================

\begin{partnote}
Format for every problem: \textbf{Problem} $\rightarrow$ \textbf{Hints}
$\rightarrow$ \textsc{attempt before reading on} $\rightarrow$
\textbf{Solution} with commentary $\rightarrow$ \textbf{Interviewer follow-up
variants}. Every solution compiles clean under
\code{gcc -Wall -Wextra -std=c11}.
\end{partnote}

%---------------------------------------------------------------
\begin{gymproblem}{Swap two integers through pointers}

\textbf{Problem.} Write \code{void swap(int *a, int *b)} that exchanges the
two values. Then explain why the same thing written as
\code{void swap(int a, int b)} cannot work.

\textbf{Hints.}
\begin{itemize}
\item You need the \emph{addresses} of the caller's variables, not copies.
\item A temporary is fine; no XOR tricks required (and XOR-swap breaks when
  both pointers are equal).
\end{itemize}

\attempt

\textbf{Solution.}
\begin{cgym}
#include <stdio.h>

void swap(int *a, int *b) {   /* a, b hold the addresses of the caller's ints */
    int tmp = *a;             /* read the value at a into a temporary          */
    *a = *b;                  /* copy b's value into a's location              */
    *b = tmp;                 /* copy the saved value into b's location        */
}

int main(void) {
    int x = 3, y = 5;
    swap(&x, &y);             /* pass addresses so swap reaches the originals  */
    printf("%d %d\n", x, y);  /* prints: 5 3                                   */
    return 0;
}
\end{cgym}

C is pass-by-value: a plain \code{swap(int a, int b)} receives \emph{copies},
swaps the copies, and the caller's variables are untouched. Passing pointers
gives the function reach into the caller's storage.

\followups
\begin{enumerate}
\item \emph{``Make it generic for any type.''} $\rightarrow$
  \code{void swap(void *a, void *b, size\_t n)} using a byte loop or
  \code{memcpy} through a temporary buffer.
\item \emph{``Why is XOR-swap dangerous?''} $\rightarrow$ If \code{a == b},
  XOR-swap zeroes the value; it also defeats the optimiser and is no faster
  on modern cores.
\item \emph{``What if a caller passes \code{swap(\&x, \&x)}?''} $\rightarrow$
  The temp-based version is safe (writes the same value back); discuss
  aliasing.
\end{enumerate}
\end{gymproblem}

%---------------------------------------------------------------
\begin{gymproblem}{\code{my\_strlen} from scratch}

\textbf{Problem.} Implement \code{size\_t my\_strlen(const char *s)} without
calling library functions.

\textbf{Hints.}
\begin{itemize}
\item Walk until the NUL terminator.
\item \code{const} because you don't modify the string. Return
  \code{size\_t}.
\end{itemize}

\attempt

\textbf{Solution.}
\begin{cgym}
#include <stddef.h>

size_t my_strlen(const char *s) {
    const char *p = s;        /* second cursor so we can subtract at the end    */
    while (*p != '\0') {      /* stop at the NUL terminator                     */
        p++;                  /* advance one char (chars are 1 byte by defn)    */
    }
    return (size_t)(p - s);   /* element distance = number of chars before NUL  */
}
\end{cgym}

\code{p - s} is pointer subtraction over \code{char}, so it yields the count
directly. The cast to \code{size\_t} is because pointer difference is the
signed \code{ptrdiff\_t}.

\followups
\begin{enumerate}
\item \emph{``What if \code{s} is NULL?''} $\rightarrow$ Standard
  \code{strlen} is UB on NULL; decide a contract (assert, or document
  non-NULL).
\item \emph{``Make it count up to a maximum (\code{strnlen}).''}
  $\rightarrow$ Add \code{size\_t maxlen} and a second loop condition.
\item \emph{``Where does this go wrong on a non-terminated buffer?''}
  $\rightarrow$ It reads past the end --- out-of-bounds; tie to why
  fixed-width fields in a packet are safer than NUL-terminated ones.
\end{enumerate}
\end{gymproblem}

%---------------------------------------------------------------
\begin{gymproblem}{\code{my\_memcpy} and the overlap question}

\textbf{Problem.} Implement
\code{void *my\_memcpy(void *dst, const void *src, size\_t n)}. Then state
what happens if the regions overlap and what you'd use instead.

\textbf{Hints.}
\begin{itemize}
\item Copy byte by byte through \code{unsigned char *}.
\item Return \code{dst} to match the standard signature.
\end{itemize}

\attempt

\textbf{Solution.}
\begin{cgym}
#include <stddef.h>

void *my_memcpy(void *dst, const void *src, size_t n) {
    unsigned char *d = dst;          /* byte cursor into destination          */
    const unsigned char *s = src;    /* byte cursor into source (read-only)   */
    while (n-- > 0) {                /* copy exactly n bytes                  */
        *d++ = *s++;                 /* post-increment both cursors           */
    }
    return dst;                      /* standard memcpy returns the dest ptr  */
}
\end{cgym}

\code{unsigned char} is the correct lens for raw bytes (well-defined, no
padding traps). If \code{dst} and \code{src} overlap, this forward copy can
clobber source bytes before reading them --- that case is \textbf{undefined
behaviour} for \code{memcpy} and is exactly why \code{memcpy} declares its
pointers \code{restrict}. Use \code{memmove}, which detects direction and
copies backward when needed.

\followups
\begin{enumerate}
\item \emph{``Write \code{my\_memmove}.''} $\rightarrow$ If \code{d < s} copy
  forward, else copy from the top down.
\item \emph{``Speed it up.''} $\rightarrow$ Copy word-at-a-time once both
  pointers are aligned, byte-tail the remainder; mention \code{restrict}.
\item \emph{``Why \code{unsigned char} and not \code{char}?''} $\rightarrow$
  \code{char} signedness is implementation-defined; \code{unsigned char} is
  the canonical byte type and is exempt from strict aliasing.
\end{enumerate}
\end{gymproblem}

%---------------------------------------------------------------
\begin{gymproblem}{Count set bits (population count)}

\textbf{Problem.} Write \code{int popcount(uint32\_t v)} returning the number
of 1-bits. Give a simple version and a faster one.

\textbf{Hints.}
\begin{itemize}
\item Naive: test bit 0, shift right, repeat 32 times.
\item Faster: \code{v \&= (v - 1)} clears the lowest set bit each iteration
  $\rightarrow$ loops only as many times as there are set bits.
\end{itemize}

\attempt

\textbf{Solution.}
\begin{cgym}
#include <stdint.h>

int popcount(uint32_t v) {
    int count = 0;
    while (v != 0u) {        /* Kernighan: loop once per set bit     */
        v &= (v - 1u);       /* clears the least-significant set bit */
        count++;
    }
    return count;
}
\end{cgym}

\code{v - 1} flips the lowest set bit to 0 and sets all lower bits to 1;
ANDing with \code{v} removes exactly that lowest set bit. So the loop runs
$k$ times for $k$ set bits, beating the fixed 32-iteration shift version on
sparse words.

\followups
\begin{enumerate}
\item \emph{``Why not just use \code{\_\_builtin\_popcount}?''} $\rightarrow$
  Fine on a host/GCC; on a bare MCU it may expand to a library call --- know
  the portable fallback.
\item \emph{``Count bits in a register status word to detect stuck bits.''}
  $\rightarrow$ Ties directly to reading a peripheral status/fault register
  during board bring-up.
\item \emph{``Parallel / bit-twiddling O(1) version?''} $\rightarrow$ The
  SWAR mask-and-add sequence; mention it exists, derive only if pushed.
\end{enumerate}
\end{gymproblem}

%---------------------------------------------------------------
\begin{gymproblem}{Set / clear / toggle / test a bit (register hygiene)}

\textbf{Problem.} Write the four canonical bit macros ---
\code{SET\_BIT}, \code{CLEAR\_BIT}, \code{TOGGLE\_BIT}, \code{TEST\_BIT} ---
and demonstrate them on a 32-bit ``register''. These four lines are the bread
and butter of every driver you will ever touch.

\textbf{Hints.}
\begin{itemize}
\item Build everything from one \code{BIT(n)} mask macro.
\item Use \code{1u}, not \code{1}: shifting a signed 1 into bit 31 is UB.
\item Parenthesise every macro argument.
\end{itemize}

\attempt

\textbf{Solution.}
\begin{cgym}
#include <stdint.h>
#include <stdio.h>

#define BIT(n)            (1u << (n))
#define SET_BIT(r, n)     ((r) |=  BIT(n))
#define CLEAR_BIT(r, n)   ((r) &= ~BIT(n))
#define TOGGLE_BIT(r, n)  ((r) ^=  BIT(n))
#define TEST_BIT(r, n)    (((r) >> (n)) & 1u)

int main(void) {
    uint32_t reg = 0u;
    SET_BIT(reg, 3);
    SET_BIT(reg, 7);
    CLEAR_BIT(reg, 3);
    TOGGLE_BIT(reg, 0);
    printf("reg = 0x%08X, bit7 = %u\n", reg, TEST_BIT(reg, 7));
    /* prints: reg = 0x00000081, bit7 = 1 */
    return 0;
}
\end{cgym}

OR sets without disturbing neighbours; AND with the complement clears; XOR
flips; shift-and-mask reads. The \code{1u} matters: \code{1 << 31} on a
32-bit \code{int} is undefined behaviour, \code{1u << 31} is fine.

\followups
\begin{enumerate}
\item \emph{``Set bits 4--7 to the value 0b1010 in one expression.''}
  $\rightarrow$ read-modify-write with a field mask:
  \code{reg = (reg \& \textasciitilde(0xFu << 4)) | (0xAu << 4);}
\item \emph{``Why can a read-modify-write on a live hardware register still
  be wrong?''} $\rightarrow$ An interrupt between the read and the write
  loses one side's update --- bridge to atomic set/clear registers and
  critical sections (Part 2 / Chapter 4).
\item \emph{``Why must \code{r} not have side effects
  (e.g.\ \code{SET\_BIT(*p++, 3)})?''} $\rightarrow$ The macro expands
  \code{r} twice; argue for an inline function or single-evaluation
  discipline.
\end{enumerate}
\end{gymproblem}

%---------------------------------------------------------------
\begin{gymproblem}{Reverse a string in place}

\textbf{Problem.} Write \code{void reverse(char *s)} that reverses a
NUL-terminated string in place.

\textbf{Hints.}
\begin{itemize}
\item Two indices/pointers from both ends moving toward the middle.
\item You need the length first (or a two-pointer walk).
\end{itemize}

\attempt

\textbf{Solution.}
\begin{cgym}
#include <stddef.h>

static size_t slen(const char *s) {     /* tiny local length helper */
    const char *p = s;
    while (*p) p++;
    return (size_t)(p - s);
}

void reverse(char *s) {
    if (s == NULL) return;               /* defensive: nothing to do  */
    size_t i = 0;
    size_t j = slen(s);
    if (j == 0) return;                  /* empty string              */
    j--;                                 /* j = index of last char    */
    while (i < j) {                      /* swap ends, walk inward    */
        char tmp = s[i];
        s[i] = s[j];
        s[j] = tmp;
        i++;
        j--;
    }
}
\end{cgym}

The guard \code{if (j == 0)} matters: \code{j} is unsigned, so decrementing 0
would wrap to a huge value and the loop would run off the end --- a classic
unsigned-underflow defect.

\followups
\begin{enumerate}
\item \emph{``Reverse only the words, keeping word order.''} $\rightarrow$
  Reverse the whole string, then reverse each word --- the classic two-pass
  trick.
\item \emph{``Why must \code{s} be writable here?''} $\rightarrow$ A string
  literal would fault; the argument should come from a \code{char[]}, not
  \code{char *lit}.
\item \emph{``Reverse the bytes of a \code{uint32\_t} instead (endian
  swap).''} $\rightarrow$ Bridges into the endianness topic in Part 2.
\end{enumerate}
\end{gymproblem}

%---------------------------------------------------------------
\begin{gymproblem}{\code{my\_strcmp} --- ordering, not just equality}

\textbf{Problem.} Implement
\code{int my\_strcmp(const char *a, const char *b)} with the standard
contract: negative if \code{a < b}, zero if equal, positive if \code{a > b}.

\textbf{Hints.}
\begin{itemize}
\item Walk while characters match \emph{and} haven't hit the terminator.
\item Compare through \code{unsigned char} --- the standard requires it, and
  it keeps high-bit bytes ordered sensibly.
\end{itemize}

\attempt

\textbf{Solution.}
\begin{cgym}
#include <stdio.h>

int my_strcmp(const char *a, const char *b) {
    while (*a != '\0' && *a == *b) {    /* advance over the common prefix     */
        a++;
        b++;
    }
    /* first mismatching (or terminating) bytes decide the order */
    return (int)(unsigned char)*a - (int)(unsigned char)*b;
}

int main(void) {
    printf("%d %d %d\n",
           my_strcmp("abc", "abc"),       /* 0          */
           my_strcmp("abc", "abd") < 0,   /* 1 (a < b)  */
           my_strcmp("b", "a") > 0);      /* 1 (a > b)  */
    return 0;
}
\end{cgym}

The loop needs only \code{*a != '\textbackslash{}0'} because if \code{*a}
equals \code{*b} and \code{*a} is NUL, both ended together. The
\code{unsigned char} casts are the part interviewers probe: with plain
(possibly signed) \code{char}, a byte like \code{0xFF} would compare as
$-1$ and order before \code{'A'}, breaking sorted tables.

\followups
\begin{enumerate}
\item \emph{``Write \code{strncmp}.''} $\rightarrow$ Add a length bound and
  stop at \code{n} as well.
\item \emph{``Why is the return contract negative/zero/positive rather than
  $-1/0/1$?''} $\rightarrow$ The standard only promises sign; never test
  \code{== -1}.
\item \emph{``Case-insensitive version?''} $\rightarrow$ \code{tolower} on
  \code{unsigned char} values --- and note locale pitfalls.
\end{enumerate}
\end{gymproblem}

%---------------------------------------------------------------
\begin{gymproblem}{\code{my\_atoi} --- parsing digits off the wire}

\textbf{Problem.} Implement \code{int my\_atoi(const char *s)}: skip leading
whitespace, accept an optional sign, convert consecutive digits, stop at the
first non-digit.

\textbf{Hints.}
\begin{itemize}
\item \code{*s - '0'} converts a digit character to its value --- the digits
  are contiguous in every C character set.
\item Accumulate in a wider type if you want headroom to discuss overflow.
\end{itemize}

\attempt

\textbf{Solution.}
\begin{cgym}
int my_atoi(const char *s) {
    long val = 0;                        /* wider accumulator         */
    int sign = 1;
    while (*s == ' ' || *s == '\t') {    /* skip leading whitespace   */
        s++;
    }
    if (*s == '+' || *s == '-') {        /* optional sign             */
        if (*s == '-') {
            sign = -1;
        }
        s++;
    }
    while (*s >= '0' && *s <= '9') {     /* consume digits            */
        val = val * 10 + (*s - '0');     /* shift-and-add in base 10  */
        s++;
    }
    return (int)(sign * val);
}
\end{cgym}

\code{my\_atoi("  -123")} $\rightarrow$ $-123$;
\code{my\_atoi("+42abc")} $\rightarrow$ 42; \code{my\_atoi("0")}
$\rightarrow$ 0 (all verified). The honest caveat to volunteer: like the real
\code{atoi}, this has no overflow protection --- a long digit string
overflows silently. That admission, plus naming \code{strtol} as the
overflow-aware replacement, is exactly what a senior interviewer is fishing
for.

\followups
\begin{enumerate}
\item \emph{``Add overflow handling.''} $\rightarrow$ Check
  \code{val > (INT\_MAX - digit) / 10} before each step, clamp like
  \code{strtol}.
\item \emph{``Parse a hex field from a telemetry log line.''} $\rightarrow$
  Extend with base-16 digits; or shift by 4 per nibble.
\item \emph{``Why is \code{*s - '0'} guaranteed to work?''} $\rightarrow$ The
  C standard requires \code{'0'..'9'} contiguous; letters have no such
  guarantee.
\end{enumerate}
\end{gymproblem}

%---------------------------------------------------------------
\begin{gymproblem}{Function-pointer dispatch table (command parser)}

\textbf{Problem.} You receive single-byte commands over UART: \code{'R'}
read, \code{'W'} write, \code{'S'} status. Replace a \code{switch} with a
function-pointer dispatch table and explain the win.

\textbf{Hints.}
\begin{itemize}
\item An array of \code{\{command, handler\}} pairs.
\item Each handler shares one signature.
\end{itemize}

\attempt

\textbf{Solution.}
\begin{cgym}
#include <stdio.h>
#include <stddef.h>

typedef int (*handler_fn)(int arg);     /* common handler signature   */

static int do_read(int arg)   { return arg + 1; }  /* stand-ins       */
static int do_write(int arg)  { return arg - 1; }
static int do_status(int arg) { return arg * 2; }

typedef struct {
    char        cmd;          /* the command byte from the wire        */
    handler_fn  fn;           /* routine that handles it               */
} command_t;

static const command_t table[] = {      /* the dispatch table          */
    { 'R', do_read   },
    { 'W', do_write  },
    { 'S', do_status },
};

static handler_fn lookup(char cmd) {
    for (size_t i = 0; i < sizeof(table) / sizeof(table[0]); i++) {
        if (table[i].cmd == cmd) {
            return table[i].fn;          /* found the handler          */
        }
    }
    return NULL;                         /* unknown command            */
}

int main(void) {
    char incoming = 'W';
    handler_fn h = lookup(incoming);
    if (h != NULL) {
        printf("result = %d\n", h(10));  /* calls do_write(10) -> 9    */
    } else {
        printf("unknown command\n");
    }
    return 0;
}
\end{cgym}

The win: adding a command is one table row, not another \code{case} plus the
risk of a missing \code{break}. The dispatch is data-driven, so the same
engine handles any protocol just by swapping the table --- which is how a
clean RS422/CAN message parser is structured. The interrupt vector table is
the same idea baked into hardware: an array of function pointers indexed by
exception number.

\followups
\begin{enumerate}
\item \emph{``What's the danger of an uninitialised function pointer?''}
  $\rightarrow$ Calling it jumps to garbage $\rightarrow$ hard fault; always
  init to NULL and check.
\item \emph{``Make it O(1) instead of O(n).''} $\rightarrow$ Index by the
  command byte directly into a 256-entry sparse table when the command space
  is small/dense.
\item \emph{``Where do function pointers appear in firmware you didn't
  write?''} $\rightarrow$ ISR vector table, RTOS task entry points, callback
  registration in driver APIs.
\end{enumerate}
\end{gymproblem}

%---------------------------------------------------------------
\begin{gymproblem}{Prove structure padding and shrink it}

\textbf{Problem.} Without running it on the target, predict \code{sizeof} for
the two structs below on a 32-bit-aligned platform, then write a program that
prints the sizes and member offsets to confirm.

\begin{cgym}
struct bad  { uint8_t a; uint32_t b; uint16_t c; };  /* poor order      */
struct good { uint32_t b; uint16_t c; uint8_t a; };  /* ordered by size */
\end{cgym}

\textbf{Hints.}
\begin{itemize}
\item \code{offsetof} from \code{<stddef.h>} reveals where each member lands.
\item Tail padding rounds the size up to the largest member's alignment.
\end{itemize}

\attempt

\textbf{Solution.}
\begin{cgym}
#include <stdio.h>
#include <stddef.h>
#include <stdint.h>

struct bad  { uint8_t a; uint32_t b; uint16_t c; };
struct good { uint32_t b; uint16_t c; uint8_t a; };

int main(void) {
    printf("bad : size=%zu  a@%zu b@%zu c@%zu\n",
           sizeof(struct bad),
           offsetof(struct bad, a),
           offsetof(struct bad, b),
           offsetof(struct bad, c));
    printf("good: size=%zu  b@%zu c@%zu a@%zu\n",
           sizeof(struct good),
           offsetof(struct good, b),
           offsetof(struct good, c),
           offsetof(struct good, a));
    return 0;
}
\end{cgym}

Predicted on a typical 32-bit-aligned ABI: \code{bad} $\rightarrow$
\code{a@0}, 3 pad, \code{b@4}, \code{c@8}, 2 tail pad $\rightarrow$
\textbf{size 12}. \code{good} $\rightarrow$ \code{b@0}, \code{c@4},
\code{a@6}, 1 tail pad $\rightarrow$ \textbf{size 8}. Same data, 33\%
smaller, just by ordering members large-to-small. (Exact numbers are
ABI-dependent; the program confirms them on whatever you run it on ---
verified 12/8 on x86-64 GCC.)

\followups
\begin{enumerate}
\item \emph{``Now make \code{bad} exactly 7 bytes.''} $\rightarrow$
  \code{\_\_attribute\_\_((packed))} --- but warn about unaligned access
  cost/fault on some ARM cores.
\item \emph{``Why does the tail get padded?''} $\rightarrow$ So
  \code{struct bad arr[2]} keeps \code{arr[1].b} aligned.
\item \emph{``How does this interact with sending the struct over RS422?''}
  $\rightarrow$ Both ends must agree on padding \emph{and} byte order, or you
  parse field-by-field instead.
\end{enumerate}
\end{gymproblem}

%---------------------------------------------------------------
\begin{gymproblem}{Type-pun a float safely (endianness probe warm-up)}

\textbf{Problem.} Print the four raw bytes of \code{float 1.0f} in memory
order, using a \emph{defined} technique (no strict-aliasing violation).

\textbf{Hints.}
\begin{itemize}
\item \code{union} member read, or \code{memcpy} into a byte array --- both
  are legal.
\item A \code{*(uint32\_t*)\&f} cast is the trap; avoid it.
\end{itemize}

\attempt

\textbf{Solution.}
\begin{cgym}
#include <stdio.h>
#include <stdint.h>
#include <string.h>

int main(void) {
    float f = 1.0f;
    unsigned char bytes[sizeof f];
    memcpy(bytes, &f, sizeof f);   /* defined: copies the object bytes */
    for (size_t i = 0; i < sizeof f; i++) {
        printf("%02X ", bytes[i]); /* little-endian: 00 00 80 3F       */
    }
    printf("\n");
    return 0;
}
\end{cgym}

\code{memcpy} into an \code{unsigned char[]} is the optimiser-safe way to
view an object's representation --- \code{unsigned char} may alias anything,
and \code{memcpy} is exempt from strict aliasing. On a little-endian host
you'll see \code{00 00 80 3F} (IEEE-754 \code{1.0f} = \code{0x3F800000}, low
byte first). On a big-endian target the order reverses --- which is the
bridge into the endianness section of Part 2 and into why telemetry floats
need explicit byte-order handling.

\followups
\begin{enumerate}
\item \emph{``Do the same with a \code{union} --- is it as legal?''}
  $\rightarrow$ Yes in C; show
  \code{union \{ float f; unsigned char b[4]; \}}.
\item \emph{``Detect endianness at runtime.''} $\rightarrow$ Pun a
  \code{uint32\_t = 1} and look at byte 0 (covered in Part 2).
\item \emph{``Why not \code{(uint32\_t)f}?''} $\rightarrow$ That
  \emph{converts} the value (truncates to 1), it does not reinterpret the
  bits --- a favourite trick question.
\end{enumerate}
\end{gymproblem}

%=====================================================================
\section{Link to Her Resume}
%=====================================================================

\begin{resumelink}
\textbf{Pointers / drivers:} ``On the BEL UART and GPIO driver verification I
worked at the register level, so I'm fluent in casting a memory-mapped
address to a \code{volatile} pointer of the correct width --- getting that
access width wrong is one of the first low-level driver defects I check
for.''

\medskip
\textbf{Structs / packets:} ``When I built the Python suite parsing
MIL-STD-1553B and RS422 telemetry, I learned the hard way to parse frames
field-by-field rather than casting a raw buffer onto a struct --- padding and
endianness differences between the sender and my parser would otherwise
corrupt the decode. The same byte-level discipline carries into the C side.''

\medskip
\textbf{Function pointers:} ``A clean message parser for CCDL or CAN traffic
is a function-pointer dispatch table, not a giant switch --- adding a message
type becomes one table entry, and it mirrors how the MCU's own interrupt
vector table works.''

\medskip
\small These connect Chapter 1 fundamentals to concrete lines already on the
resume (LPC3250/ARM9 driver work, the 1553B/RS422 Python automation, RS422
interface routines at TASL). Keep the claims to what the resume states ---
the deep project scripting happens later, only from \code{inputs/INTAKE.md}.
\end{resumelink}

%=====================================================================
\section{Self-Test Scorecard}
%=====================================================================

\begin{scorecard}
\begin{enumerate}
\item On a 32-bit MCU, what is \code{sizeof(double *)}?
\item \code{int *p; p++;} --- by how many bytes does \code{p} move?
\item Read aloud: \code{char *const p}. What is fixed, what is free?
\item Where does array-to-pointer decay \textbf{not} happen? Name two places.
\item What single byte ends a C string, and what is its value?
\item \code{arr[i]} is shorthand for what pointer expression? Is
  \code{i[arr]} legal?
\item Why is \code{char s[] = "hi"} writable but \code{char *s = "hi"} not
  safe to write?
\item Given \code{uint8\_t; uint32\_t; uint16\_t} in a struct, what is the
  size on a 32-bit ABI, and how do you shrink it?
\item Reading a union member you didn't last write --- legal in C? What about
  \code{*(uint32\_t*)\&myfloat}?
\item Name two reasons a bitfield struct must not define an on-the-wire
  protocol.
\end{enumerate}
\end{scorecard}

\begin{answerkey}
\begin{enumerate}
\item 4 bytes --- pointer size follows the target address width, not the
  pointee.
\item \code{sizeof(int)} bytes, normally 4.
\item \code{p} (the pointer) is fixed/const; \code{*p} (the pointed-to char)
  is writable.
\item Under \code{sizeof}, under unary \code{\&}, and as a \code{char[]}
  string-literal initialiser (any two).
\item The NUL byte \code{'\textbackslash{}0'}, value 0.
\item \code{*(arr + i)}; yes, \code{i[arr]} is legal because addition
  commutes.
\item The array copies the bytes into writable storage you own; the pointer
  aims at a read-only string literal --- writing it is UB.
\item 12 bytes (with 5 padding); reorder largest-to-smallest $\rightarrow$
  8 bytes. Packing is the last resort, for wire/register layouts.
\item Union member read is legal type-punning in C; the pointer-cast version
  violates strict aliasing and can be miscompiled at \code{-O2}.
\item Bit allocation order (MSB/LSB-first) is implementation-defined;
  signedness of plain-\code{int} bitfields is implementation-defined;
  straddling storage units and padding are also unspecified --- so the layout
  isn't portable across compilers/architectures.
\end{enumerate}
\end{answerkey}

\begin{partnote}
\emph{End of Chapter 1, Part 1. Part 2 (\code{volatile} / \code{const} /
\code{static} / storage classes / typecasting / preprocessor / endianness)
continues in this file in the next session.}
\end{partnote}
